\documentclass[14pt]{extarticle}
%\documentclass[a4paper]{extarticle}
%\documentclass[12pt, a4paper]{article}

% === Кодировки и язык ===
\usepackage[T2A]{fontenc}          % Кириллическая кодировка шрифтов
\usepackage[utf8]{inputenc}        % Кодировка исходного файла - UTF-8
\usepackage[russian]{babel}        % Поддержка русского языка

% === Основные пакеты ===
\usepackage{graphicx}              % Работа с изображениями
\usepackage{amsmath}               % Математические формулы
\usepackage{indentfirst}           % Отступ в первом абзаце раздела
\usepackage{fancyvrb}              % Расширенный режим Verbatim
\usepackage{fvextra}               % Поддержка breaklines и UTF-8 в Verbatim
\usepackage{comment}               % Блоки комментариев
%\usepackage{enumitem}              % Гибкая настройка списков
%\usepackage[skip=0pt]{parskip}     % Убирает лишние отступы между абзацами
\usepackage{float}
%\usepackage{makecell}

\usepackage{pdflscape}

\usepackage{listings}
% Переименовываем стандартное слово "Listing" в "Листинг"
\renewcommand{\lstlistingname}{Листинг}
\usepackage{amssymb} % Требуется для символа \hookrightarrow
\lstset{
    %basicstyle=\ttfamily,         % Моноширинный шрифт
    numbers=left,                 % Номера строк слева
    numberstyle=\tiny,            % Размер шрифта для номеров строк
    frame=single,                 % Одинарная рамка вокруг кода
    breaklines=true,              % Автоматический перенос длинных строк
    postbreak=\mbox{\tiny\ensuremath{\hookrightarrow}\space},
    %tabsize=4,                    % Размер табуляции
    %showstringspaces=false        % Отключение визуализации пробелов в строках
    }

% === Разметка страницы и вспомогательные пакеты ===
\usepackage{geometry}              % Настройка полей - используется ниже вместо ручных команд

% === Гиперссылки ===
\usepackage{hyperref}

% === Геометрия страницы ===
\geometry{
  % textheight=24cm,
  % textwidth=16cm,
  left=2.5cm,
  top=1cm
}